\begin{thebibliography}{99}\small\setlength{\itemsep}{2pt}\setlength{\parskip}{0pt}\addcontentsline{toc}{chapter}{\bibname}
\bibitem{AC} A. Adhikari and S. Cao. Correlation decay for finite lattice gauge theories at weak coupling. \emph{Annals of Probability} 53 (2025) 140.
\bibitem{ASWE} M. Aissen, I. J. Schoenberg and A. Whitney. On the generating functions of totally positive sequences I. \emph{Journal d'Analyse Math\'ematique} 2 (1952) 93.
\bibitem{Ai82} M. Aizenman. Geometric analysis of $\varphi^4$ fields and Ising models, parts I and II. \emph{Communications in Mathematical Physics} 86 (1982) 1.
\bibitem{ADC21} M. Aizenman and H. Duminil Copin. Marginal triviality of the scaling limits of critical 4D Ising and $\varphi^4_4$ models. \emph{Annals of Mathematics} 194 (2021) 163.
\bibitem{APE87} M. Albanese et al. (APE Collaboration). Glueball masses and string tension in lattice QCD. \emph{Physics Letters B} 192 (1987) 163.
\bibitem{Al86} N. Alon. Eigenvalues and expanders. \emph{Combinatorica} 6 (1986) 83.
\bibitem{Alp26} L. Alp\"oge. Announcement, with Lean code, of a proof of the Riemann hypothesis and of the generalized Riemann hypothesis for Dirichlet $L$ functions, found working with Claude, 20 November 2026, \url{https://alpo.ge/remark.pdf}. Reported in \emph{Scientific American}, 20 November 2026.
\bibitem{Am74} D. E. Amos. Computation of modified Bessel functions and their ratios. \emph{Mathematics of Computation} 28 (1974) 239.
\bibitem{ALO20} N. Anari, K. Liu and S. Oveis Gharan. Spectral independence in high dimensional expanders and applications to the hardcore model. \emph{Proceedings of FOCS}, 2020.
\bibitem{AK17} P. D. Anderson and M. Kruczenski. Loop equations and bootstrap methods in the lattice. \emph{Nuclear Physics B} 921 (2017) 702.
\bibitem{AM74} D. F. Andrews and C. L. Mallows. Scale mixtures of normal distributions. \emph{Journal of the Royal Statistical Society B} 36 (1974) 99.
\bibitem{AS12} M. Anshelevich and A. N. Sengupta. Quantum free Yang Mills on the plane. \emph{Journal of Geometry and Physics} 62 (2012) 330.
\bibitem{AK91} D. Applegate and R. Kannan. Sampling and integration of near log concave functions. \emph{Proceedings of the 23rd Annual ACM Symposium on Theory of Computing} (1991) 156.
\bibitem{As70} T. Asano. Theorems on the partition functions of the Heisenberg ferromagnets. \emph{Journal of the Physical Society of Japan} 29 (1970) 350.
\bibitem{Ba86} C. Bachas. Convexity of the quarkonium potential. \emph{Physical Review D} 33 (1986) 2723.
\bibitem{BG68} G. Backus and F. Gilbert. The resolving power of gross earth data. \emph{Geophysical Journal of the Royal Astronomical Society} 16 (1968) 169.
\bibitem{BE} D. Bakry and M. \'Emery. Diffusions hypercontractives. \emph{S\'eminaire de Probabilit\'es} XIX, Lecture Notes in Mathematics 1123, Springer, 1985.
\bibitem{BCG08} D. Bakry, P. Cattiaux and A. Guillin. Rate of convergence for ergodic continuous Markov processes: Lyapunov versus Poincar\'e. \emph{Journal of Functional Analysis} 254 (2008) 727.
\bibitem{Ba87} T. Ba\l aban. Renormalization group approach to lattice gauge field theories I and II. \emph{Communications in Mathematical Physics} 109 (1987) 249 and 116 (1988) 1.
\bibitem{Be29} S. Bernstein. Sur les fonctions absolument monotones. \emph{Acta Mathematica} 52 (1929) 1.
\bibitem{BG93} J. Besag and P. J. Green. Spatial statistics and Bayesian computation. \emph{Journal of the Royal Statistical Society B} 55 (1993) 25.
\bibitem{BYM} J. Besag, J. York and A. Molli\'e. Bayesian image restoration, with two applications in spatial statistics. \emph{Annals of the Institute of Statistical Mathematics} 43 (1991) 1.
\bibitem{Be74} J. Besag. Spatial interaction and the statistical analysis of lattice systems. \emph{Journal of the Royal Statistical Society B} 36 (1974) 192.
\bibitem{Be75} J. Besag. Statistical analysis of non lattice data. \emph{The Statistician} 24 (1975) 179.
\bibitem{BC81} G. Bhanot and M. Creutz. Variant actions and phase structure in lattice gauge theory. \emph{Physical Review D} 24 (1981) 3212.
\bibitem{BPY} P. Biane, J. Pitman and M. Yor. Probability laws related to the Jacobi theta and Riemann zeta functions, and Brownian excursions. \emph{Bulletin of the American Mathematical Society} 38 (2001) 435.
\bibitem{Bi97} P. Biane. Free Brownian motion, free stochastic calculus and random matrices. \emph{Fields Institute Communications} 12 (1997) 1.
\bibitem{BP} J. R. Birge and N. G. Polson. Optimization by slice sampling.
\bibitem{Bl47} D. Blackwell. Conditional expectation and unbiased sequential estimation. \emph{Annals of Mathematical Statistics} 18 (1947) 105.
\bibitem{BS73} P. M. Bleher and Ya. G. Sinai. Investigation of the critical point in models of the type of Dyson's hierarchical model. \emph{Communications in Mathematical Physics} 33 (1973) 23.
\bibitem{Bc33} S. Bochner. Monotone Funktionen, Stieltjessche Integrale und harmonische Analyse. \emph{Mathematische Annalen} 108 (1933) 378.
\bibitem{Bo92} L. Bondesson. \emph{Generalized Gamma Convolutions and Related Classes of Distributions and Densities}. Lecture Notes in Statistics 76, Springer, 1992.
\bibitem{BB09a} J. Borcea and P. Br\"and\'en. The Lee Yang and P\'olya Schur programs I: linear operators preserving stability. \emph{Inventiones Mathematicae} 177 (2009) 541.
\bibitem{BB09b} J. Borcea and P. Br\"and\'en. The Lee Yang and P\'olya Schur programs II: theory of stable polynomials and applications. \emph{Communications on Pure and Applied Mathematics} 62 (2009) 1595.
\bibitem{Bo96} J. Bourgain. Invariant measures for the 2D defocusing nonlinear Schr\"odinger equation. \emph{Communications in Mathematical Physics} 176 (1996) 421.
\bibitem{BC23} B. Bringmann and S. Cao. A para controlled approach to the stochastic Yang Mills equation in two dimensions. arXiv:2305.07197.
\bibitem{BC24} B. Bringmann and S. Cao. Global well posedness of the stochastic abelian Higgs equations in two dimensions. arXiv:2403.16878.
\bibitem{BCHZ} B. Bringmann, S. Cao, M. Hairer and W. Zhao. Construction of the 2D non abelian Yang Mills Higgs measure via stochastic quantisation. arXiv:2611.03640, 2026.
\bibitem{dB50} N. G. de Bruijn. The roots of trigonometric integrals. \emph{Duke Mathematical Journal} 17 (1950) 197.
\bibitem{BFS82} D. Brydges, J. Fr\"ohlich and T. Spencer. The random walk representation of classical spin systems and correlation inequalities. \emph{Communications in Mathematical Physics} 83 (1982) 123.
\bibitem{BZ03} J. Budczies and M. R. Zirnbauer. Howe duality for an induced model of lattice U($N$) Yang Mills theory. arXiv:math-ph/0305058, 2003.
\bibitem{CM82} N. Cabibbo and E. Marinari. A new method for updating SU(N) matrices in computer simulations of gauge theories. \emph{Physics Letters B} 119 (1982) 387.
\bibitem{CC24} S. Cao and S. Chatterjee. A state space for 3D Euclidean Yang Mills theories. \emph{Communications in Mathematical Physics} 405 (2024).
\bibitem{Ca20} S. Cao. Wilson loop expectations in lattice gauge theories with finite gauge groups. \emph{Communications in Mathematical Physics} 380 (2020) 1439.
\bibitem{CC23} S. Cao and S. Chatterjee. The Yang Mills heat flow with random distributional initial data. \emph{Communications in Partial Differential Equations} 48 (2023) 209.
\bibitem{CPS25} S. Cao, M. Park and S. Sheffield. Random surfaces and lattice Yang Mills. \emph{Communications of the American Mathematical Society} 5 (2025) 774.
\bibitem{CNS25} S. Cao, R. Nissim and S. Sheffield. Expanded regimes of area law for lattice Yang Mills theories. arXiv:2505.16585; Dynamical approach to area law for lattice Yang Mills. arXiv:2509.04688.
\bibitem{CPS10} C. M. Carvalho, N. G. Polson and J. G. Scott. The horseshoe estimator for sparse signals. \emph{Biometrika} 97 (2010) 465.
\bibitem{CJLP} C. M. Carvalho, M. S. Johannes, H. F. Lopes and N. G. Polson. Particle learning and smoothing. \emph{Statistical Science} 25 (2010) 88.
\bibitem{CMGC09} M. Chaleyat Maurel and V. Genon Catalot. Filtering the Wright Fisher diffusion. \emph{ESAIM Probability and Statistics} 13 (2009) 197.
\bibitem{CCHS} A. Chandra, I. Chevyrev, M. Hairer and H. Shen. Langevin dynamic for the 2D Yang Mills measure. \emph{Publications Math\'ematiques de l'IH\'ES}, 2022.
\bibitem{CCHS3} A. Chandra, I. Chevyrev, M. Hairer and H. Shen. Stochastic quantisation of Yang Mills Higgs in 3D. \emph{Inventiones Mathematicae} 237 (2024) 541.
\bibitem{Ch19so} S. Chatterjee. Rigorous solution of strongly coupled SO(N) lattice gauge theory in the large N limit. \emph{Communications in Mathematical Physics} 366 (2019) 203.
\bibitem{Ch19} S. Chatterjee. Yang Mills for probabilists. In \emph{Probability and Analysis in Interacting Physical Systems}, Springer, 2019.
\bibitem{ChSh23} I. Chevyrev and H. Shen. Invariant measure and universality of the 2D Yang Mills Langevin dynamic. arXiv:2302.12160.
\bibitem{CT62} Z. Ciesielski and S. J. Taylor. First passage times and sojourn times for Brownian motion in space and the exact Hausdorff measure of the sample path. \emph{Transactions of the American Mathematical Society} 103 (1962) 434.
\bibitem{Co92} F. Comets. On consistency of a class of estimators for exponential families of Markov random fields on the lattice. \emph{Annals of Statistics} 20 (1992) 455.
\bibitem{CMS} D. Cordero Erausquin, R. J. McCann and M. Schmuckenschl\"ager. A Riemannian interpolation inequality \`a la Borell, Brascamp and Lieb. \emph{Inventiones Mathematicae} 146 (2001) 219.
\bibitem{CT06} T. M. Cover and J. A. Thomas. \emph{Elements of Information Theory}, second edition. Wiley, 2006.
\bibitem{Cr80} M. Creutz. Monte Carlo study of quantized SU(2) gauge theory. \emph{Physical Review D} 21 (1980) 2308.
\bibitem{CJK18} D. Crisan, C. Janjigian and T. G. Kurtz. Particle representations for stochastic partial differential equations with boundary conditions. \emph{Electronic Journal of Probability} 23 (2018).
\bibitem{DWW99} P. Damien, J. Wakefield and S. Walker. Gibbs sampling for Bayesian non conjugate and hierarchical models by using auxiliary variables. \emph{Journal of the Royal Statistical Society B} 61 (1999) 331.
\bibitem{De74} P. Deligne. La conjecture de Weil I. \emph{Publications Math\'ematiques de l'IH\'ES} 43 (1974) 273.
\bibitem{DDLMZ} A. Denbleyker, D. Du, Y. Liu, Y. Meurice and H. Zou. Fisher's zeros as the boundary of renormalization group flows in complex coupling spaces. \emph{Physical Review Letters} 104 (2010) 251601.
\bibitem{DT83} D. M. DeTurck. Deforming metrics in the direction of their Ricci tensors. \emph{Journal of Differential Geometry} 18 (1983) 157.
\bibitem{DS81} P. Diaconis and M. Shahshahani. Generating a random permutation with random transpositions. \emph{Zeitschrift f\"ur Wahrscheinlichkeitstheorie} 57 (1981) 159.
\bibitem{DiZa92} I. H. Dinwoodie and S. L. Zabell. Large deviations for exchangeable random vectors. \emph{Annals of Probability} 20 (1992) 1147.
\bibitem{DLCL} Dirichlet L functions and the critical line. Unpublished expanded account of the argument announced by L. Alp\"oge on 20 November 2026, written to accompany its Lean development; 55 pages, no author named. Copy seen by the authors, November 2026.
\bibitem{DS85} R. L. Dobrushin and S. B. Shlosman. Completely analytical Gibbs fields. In \emph{Statistical Physics and Dynamical Systems}, Birkh\"auser, 1985, 371.
\bibitem{DS87} R. L. Dobrushin and S. B. Shlosman. Completely analytical interactions: constructive description. \emph{Journal of Statistical Physics} 46 (1987) 983.
\bibitem{DV75} M. D. Donsker and S. R. S. Varadhan. Asymptotic evaluation of certain Markov process expectations for large time I. \emph{Communications on Pure and Applied Mathematics} 28 (1975) 1.
\bibitem{DK} M. R. Douglas and V. A. Kazakov. Large N phase transition in continuum QCD in two dimensions. \emph{Physics Letters B} 319 (1993) 219.
\bibitem{Dr89} B. K. Driver. YM2: continuum expectations, lattice convergence, and lassos. \emph{Communications in Mathematical Physics} 123 (1989) 575.
\bibitem{DZ} J. M. Drouffe and J. B. Zuber. Strong coupling and mean field methods in lattice gauge theories. \emph{Physics Reports} 102 (1983) 1.
\bibitem{DKPR87} S. Duane, A. D. Kennedy, B. J. Pendleton and D. Roweth. Hybrid Monte Carlo. \emph{Physics Letters B} 195 (1987) 216.
\bibitem{Dy69} F. J. Dyson. Existence of a phase transition in a one dimensional Ising ferromagnet. \emph{Communications in Mathematical Physics} 12 (1969) 91.
\bibitem{Ed52} A. Edrei. On the generating functions of totally positive sequences II. \emph{Journal d'Analyse Math\'ematique} 2 (1952) 104.
\bibitem{ES88} R. G. Edwards and A. D. Sokal. Generalization of the Fortuin Kasteleyn Swendsen Wang representation and Monte Carlo algorithm. \emph{Physical Review D} 38 (1988) 2009.
\bibitem{El75} S. Elitzur. Impossibility of spontaneously breaking local symmetries. \emph{Physical Review D} 12 (1975) 3978.
\bibitem{Ew72} W. J. Ewens. The sampling theory of selectively neutral alleles. \emph{Theoretical Population Biology} 3 (1972) 87.
\bibitem{FP67} L. D. Faddeev and V. N. Popov. Feynman diagrams for the Yang Mills field. \emph{Physics Letters B} 25 (1967) 29.
\bibitem{FGGS00} E. Farhi, J. Goldstone, S. Gutmann and M. Sipser. Quantum computation by adiabatic evolution. arXiv:quant-ph/0001106, 2000.
\bibitem{Fav22} M. Favoni, A. Ipp, D. I. M\"uller and D. Schuh. Lattice gauge equivariant convolutional neural networks. \emph{Physical Review Letters} 128 (2022) 032003.
\bibitem{FFS92} R. Fern\'andez, J. Fr\"ohlich and A. D. Sokal. \emph{Random Walks, Critical Phenomena, and Triviality in Quantum Field Theory}. Springer, 1992.
\bibitem{FHK} J. Fingberg, U. Heller and F. Karsch. Scaling and asymptotic scaling in the SU(2) gauge theory. \emph{Nuclear Physics B} 392 (1993) 493.
\bibitem{FLV22} M. P. Forsstr\"om, J. Lenells and F. Viklund. Wilson loops in finite abelian lattice gauge theories. \emph{Annales de l'Institut Henri Poincar\'e Probabilit\'es et Statistiques} 58 (2022) 2129.
\bibitem{FK72} C. M. Fortuin and P. W. Kasteleyn. On the random cluster model I: introduction and relation to other models. \emph{Physica} 57 (1972) 536.
\bibitem{FKG71} C. M. Fortuin, P. W. Kasteleyn and J. Ginibre. Correlation inequalities on some partially ordered sets. \emph{Communications in Mathematical Physics} 22 (1971) 89.
\bibitem{FP08} N. Friel and A. N. Pettitt. Marginal likelihood estimation via power posteriors. \emph{Journal of the Royal Statistical Society B} 70 (2008) 589.
\bibitem{FKP} A. Frieze, R. Kannan and N. G. Polson. Sampling from log concave distributions. \emph{Annals of Applied Probability} 4 (1994) 812.
\bibitem{Fr82} J. Fr\"ohlich. On the triviality of $\lambda\varphi^4_d$ theories and the approach to the critical point in $d \ge 4$ dimensions. \emph{Nuclear Physics B} 200 (1982) 281.
\bibitem{FS82} J. Fr\"ohlich and T. Spencer. Massless phases and symmetry restoration in abelian gauge theories and spin systems. \emph{Communications in Mathematical Physics} 83 (1982) 411.
\bibitem{GKSW} D. Gaiotto, A. Kapustin, N. Seiberg and B. Willett. Generalized global symmetries. \emph{Journal of High Energy Physics} 02 (2015) 172.
\bibitem{GSV16} A. Galanis, D. \v{S}tefankovi\v{c} and E. Vigoda. Inapproximability for antiferromagnetic spin systems in the tree nonuniqueness region. \emph{Journal of the ACM} 62 (2015) article 50.
\bibitem{Gar73} H. Garland. $p$ adic curvature and the cohomology of discrete subgroups of $p$ adic groups. \emph{Annals of Mathematics} 97 (1973) 375.
\bibitem{Gi70} J. Ginibre. General formulation of Griffiths' inequalities. \emph{Communications in Mathematical Physics} 16 (1970) 310.
\bibitem{GJ} J. Glimm and A. Jaffe. \emph{Quantum Physics: A Functional Integral Point of View}. Springer, second edition, 1987.
\bibitem{Go67} C. M. Goldie. A class of infinitely divisible random variables. \emph{Proceedings of the Cambridge Philosophical Society} 63 (1967) 1141.
\bibitem{GS89} J. Goodman and A. D. Sokal. Multigrid Monte Carlo method: conceptual foundations. \emph{Physical Review D} 40 (1989) 2035.
\bibitem{GM82} M. G\"opfert and G. Mack. Proof of confinement of static quarks in three dimensional U(1) lattice gauge theory for all values of the coupling constant. \emph{Communications in Mathematical Physics} 82 (1982) 545.
\bibitem{Gr79} J. P. Greensite. Calculation of the Yang Mills vacuum wave functional. \emph{Nuclear Physics B} 158 (1979) 469.
\bibitem{Gr78} V. N. Gribov. Quantization of non abelian gauge theories. \emph{Nuclear Physics B} 139 (1978) 1.
\bibitem{Gri67} R. B. Griffiths. Correlations in Ising ferromagnets I. \emph{Journal of Mathematical Physics} 8 (1967) 478.
\bibitem{Gri80} R. C. Griffiths. Lines of descent in the diffusion approximation of neutral Wright Fisher models. \emph{Theoretical Population Biology} 17 (1980) 37.
\bibitem{GW80} D. J. Gross and E. Witten. Possible third order phase transition in the large N lattice gauge theory. \emph{Physical Review D} 21 (1980) 446.
\bibitem{GW73} D. J. Gross and F. Wilczek. Ultraviolet behavior of nonabelian gauge theories. \emph{Physical Review Letters} 30 (1973) 1343.
\bibitem{GKS89} L. Gross, C. King and A. Sengupta. Two dimensional Yang Mills theory via stochastic differential equations. \emph{Annals of Physics} 194 (1989) 65.
\bibitem{Gro75} L. Gross. Logarithmic Sobolev inequalities. \emph{American Journal of Mathematics} 97 (1975) 1061.
\bibitem{GRS75} F. Guerra, L. Rosen and B. Simon. The $P(\phi)_2$ Euclidean quantum field theory as classical statistical mechanics. \emph{Annals of Mathematics} 101 (1975) 111.
\bibitem{Gu80} A. H. Guth. Existence proof of a nonconfining phase in four dimensional U(1) lattice gauge theory. \emph{Physical Review D} 21 (1980) 2291.
\bibitem{Ha14} M. Hairer. A theory of regularity structures. \emph{Inventiones Mathematicae} 198 (2014) 269.
\bibitem{Ha93} R. S. Hamilton. A matrix Harnack estimate for the heat equation. \emph{Communications in Analysis and Geometry} 1 (1993) 113.
\bibitem{Ha93b} R. S. Hamilton. Monotonicity formulas for parabolic flows on manifolds. \emph{Communications in Analysis and Geometry} 1 (1993) 127.
\bibitem{HLT} M. Hansen, A. Lupo and N. Tantalo. Extraction of spectral densities from lattice correlators. \emph{Physical Review D} 99 (2019) 094508.
\bibitem{HN75} G. Harder and M. S. Narasimhan. On the cohomology groups of moduli spaces of vector bundles on curves. \emph{Mathematische Annalen} 212 (1975) 215.
\bibitem{HW74} P. Hartman and G. S. Watson. Normal distribution functions on spheres and the modified Bessel functions. \emph{Annals of Probability} 2 (1974) 593.
\bibitem{Har76} P. Hartman. Completely monotone families of solutions of $n$th order linear differential equations and infinitely divisible distributions. \emph{Annali della Scuola Normale Superiore di Pisa, Classe di Scienze} (4) 3 (1976) 267.
\bibitem{HL72} O. J. Heilmann and E. H. Lieb. Theory of monomer dimer systems. \emph{Communications in Mathematical Physics} 25 (1972) 190.
\bibitem{HS95} J. Hesselbo and R. B. Stinchcombe. Monte Carlo simulation and global optimization without parameters. \emph{Physical Review Letters} 74 (1995) 2151.
\bibitem{Hi98} D. Higdon. Auxiliary variable methods for Markov chain Monte Carlo with applications. \emph{Journal of the American Statistical Association} 93 (1998) 585.
\bibitem{tH78} G. 't Hooft. On the phase transition towards permanent quark confinement. \emph{Nuclear Physics B} 138 (1978) 1.
\bibitem{tHV72} G. 't Hooft and M. Veltman. Regularization and renormalization of gauge fields. \emph{Nuclear Physics B} 44 (1972) 189.
\bibitem{tH74} G. 't Hooft. A planar diagram theory for strong interactions. \emph{Nuclear Physics B} 72 (1974) 461.
\bibitem{tH71} G. 't Hooft. Renormalizable Lagrangians for massive Yang Mills fields. \emph{Nuclear Physics B} 35 (1971) 167.
\bibitem{tH71a} G. 't Hooft. Renormalization of massless Yang Mills fields. \emph{Nuclear Physics B} 33 (1971) 173.
\bibitem{Ho89} R. Howe. Remarks on classical invariant theory. \emph{Transactions of the American Mathematical Society} 313 (1989) 539.
\bibitem{IS} K. R. Ito and E. Seiler. On the recent paper on quark confinement by Tomboulis. arXiv:0711.4930, 2007; Further discussion of Tomboulis' approach to the confinement problem. arXiv:0803.3019, 2008.
\bibitem{Ito85} K. R. Ito. Permanent quark confinement in 4D hierarchical lattice gauge theories of Migdal Kadanoff type. \emph{Physical Review Letters} 55 (1985) 558.
\bibitem{JW} A. Jaffe and E. Witten. Quantum Yang Mills theory. In \emph{The Millennium Prize Problems}, Clay Mathematics Institute, 2006.
\bibitem{JS93} M. Jerrum and A. Sinclair. Polynomial time approximation algorithms for the Ising model. \emph{SIAM Journal on Computing} 22 (1993) 1087.
\bibitem{JPS09} M. S. Johannes, N. G. Polson and J. R. Stroud. Optimal filtering of jump diffusions: extracting latent states from asset prices. \emph{Review of Financial Studies} 22 (2009) 2559.
\bibitem{JLP12} S. P. Jordan, K. S. M. Lee and J. Preskill. Quantum algorithms for quantum field theories. \emph{Science} 336 (2012) 1130.
\bibitem{Ka49} M. Kac. On distributions of certain Wiener functionals. \emph{Transactions of the American Mathematical Society} 65 (1949) 1.
\bibitem{Ka66} M. Kac. Can one hear the shape of a drum? \emph{American Mathematical Monthly} 73 (1966) 1.
\bibitem{Ka74} M. Kac. Comments on [93]. In G. P\'olya, \emph{Collected Papers, Volume II: Location of Zeros}, edited by R. P. Boas, MIT Press, 1974, pages 424 to 426.
\bibitem{Ka69} M. Kac. Some mathematical models in science. \emph{Science} 166 (1969) 695.
\bibitem{Ka76} L. P. Kadanoff. Notes on Migdal's recursion formulas. \emph{Annals of Physics} 100 (1976) 359.
\bibitem{Ka75} L. P. Kadanoff. Variational principles and approximate renormalization group calculations. \emph{Physical Review Letters} 34 (1975) 1005.
\bibitem{Kan20} G. Kanwar and coauthors. Equivariant flow based sampling for lattice gauge theory. \emph{Physical Review Letters} 125 (2020) 121601.
\bibitem{KKN} D. Karabali, C. Kim and V. P. Nair. On the vacuum wavefunction and string tension of Yang Mills theories in (2+1) dimensions. \emph{Physics Letters B} 434 (1998) 103.
\bibitem{KM59} S. Karlin and J. McGregor. Coincidence probabilities. \emph{Pacific Journal of Mathematics} 9 (1959) 1141.
\bibitem{Ka68} S. Karlin. \emph{Total Positivity}. Stanford University Press, 1968.
\bibitem{Kau49} B. Kaufman. Crystal statistics II. Partition function evaluated by spinor analysis. \emph{Physical Review} 76 (1949) 1232.
\bibitem{KZ23} V. Kazakov and Z. Zheng. Bootstrap for lattice Yang Mills theory. \emph{Physical Review D} 107 (2023) L051501.
\bibitem{Ke79} J. Keilson. \emph{Markov Chain Models: Rarity and Exponentiality}. Springer, 1979.
\bibitem{KP85} A. D. Kennedy and B. J. Pendleton. Improved heatbath method for Monte Carlo calculations in lattice gauge theories. \emph{Physics Letters B} 156 (1985) 393.
\bibitem{Ke78} J. Kent. Some probabilistic properties of Bessel functions. \emph{Annals of Probability} 6 (1978) 760.
\bibitem{Kes59} H. Kesten. Symmetric random walks on groups. \emph{Transactions of the American Mathematical Society} 92 (1959) 336.
\bibitem{KS03} H. H. Kim and P. Sarnak. Refined estimates towards the Ramanujan and Selberg conjectures. Appendix to H. H. Kim, \emph{Journal of the American Mathematical Society} 16 (2003) 139.
\bibitem{Ki78} J. F. C. Kingman. The representation of partition structures. \emph{Journal of the London Mathematical Society} 18 (1978) 374.
\bibitem{Ki63} J. F. C. Kingman. Random walks with spherical symmetry. \emph{Acta Mathematica} 109 (1963) 11.
\bibitem{Ki73} J. F. C. Kingman. Subadditive ergodic theory. \emph{Annals of Probability} 1 (1973) 883.
\bibitem{Ki82} J. F. C. Kingman. The coalescent. \emph{Stochastic Processes and their Applications} 13 (1982) 235.
\bibitem{KV86} C. Kipnis and S. R. S. Varadhan. Central limit theorem for additive functionals of reversible Markov processes and applications to simple exclusions. \emph{Communications in Mathematical Physics} 104 (1986) 1.
\bibitem{KS75} J. Kogut and L. Susskind. Hamiltonian formulation of Wilson's lattice gauge theories. \emph{Physical Review D} 11 (1975) 395.
\bibitem{Ko36} B. O. Koopman. On distributions admitting a sufficient statistic. \emph{Transactions of the American Mathematical Society} 39 (1936) 399.
\bibitem{KN97} S. M. Krone and C. Neuhauser. Ancestral processes with selection. \emph{Theoretical Population Biology} 51 (1997) 210.
\bibitem{KX04} T. G. Kurtz and J. Xiong. A stochastic evolution equation arising from the fluctuations of a class of interacting particle systems. \emph{Communications in Mathematical Sciences} 2 (2004) 325.
\bibitem{KX01} T. G. Kurtz and J. Xiong. Numerical solutions for a class of SPDEs with application to filtering. In \emph{Stochastics in Finite and Infinite Dimensions}, Birkh\"auser, 2001, 233.
\bibitem{KX99} T. G. Kurtz and J. Xiong. Particle representations for a class of nonlinear SPDEs. \emph{Stochastic Processes and their Applications} 83 (1999) 103.
\bibitem{Ku64} H. J. Kushner. On the differential equations satisfied by conditional probability densities of Markov processes, with applications. \emph{SIAM Journal on Control} 2 (1964) 106.
\bibitem{LLR} K. Langfeld, B. Lucini and A. Rago. Density of states in gauge theories. \emph{Physical Review Letters} 109 (2012) 111601.
\bibitem{LH74} C. L. Lawson and R. J. Hanson. \emph{Solving Least Squares Problems}. Prentice Hall, 1974 (SIAM reprint, 1995).
\bibitem{LY52} T. D. Lee and C. N. Yang. Statistical theory of equations of state and phase transitions II. \emph{Physical Review} 87 (1952) 410.
\bibitem{Le03} T. L\'evy. Yang Mills measure on compact surfaces. \emph{Memoirs of the American Mathematical Society} 166 (2003).
\bibitem{Le17} T. L\'evy. The master field on the plane. \emph{Ast\'erisque} 388 (2017).
\bibitem{LSV05} A. Lubotzky, B. Samuels and U. Vishne. Ramanujan complexes of type $\tilde A_d$. \emph{Israel Journal of Mathematics} 149 (2005) 267.
\bibitem{LPS88} A. Lubotzky, R. Phillips and P. Sarnak. Ramanujan graphs. \emph{Combinatorica} 8 (1988) 261.
\bibitem{LS11} M. L\"uscher and S. Schaefer. Lattice QCD without topology barriers. \emph{Journal of High Energy Physics} 07 (2011) 036.
\bibitem{LW01} M. L\"uscher and P. Weisz. Locality and exponential error reduction in numerical lattice gauge theory. \emph{Journal of High Energy Physics} 09 (2001) 010.
\bibitem{Lu77} M. L\"uscher. Construction of a selfadjoint, strictly positive transfer matrix for Euclidean lattice gauge theories. \emph{Communications in Mathematical Physics} 54 (1977) 283.
\bibitem{Lu10} M. L\"uscher. Properties and uses of the Wilson flow in lattice QCD. \emph{Journal of High Energy Physics} 08 (2010) 071.
\bibitem{MP80} G. Mack and V. B. Petkova. Sufficient condition for confinement of static quarks by a vortex condensation mechanism. \emph{Annals of Physics} 125 (1980) 117.
\bibitem{MS88} N. Madras and A. D. Sokal. The pivot algorithm: a highly efficient Monte Carlo method for the self avoiding walk. \emph{Journal of Statistical Physics} 50 (1988) 109.
\bibitem{MRS} J. Magnen, V. Rivasseau and R. S\'eneor. Construction of YM$_4$ with an infrared cutoff. \emph{Communications in Mathematical Physics} 155 (1993) 325.
\bibitem{Malo} E. Malo. Note sur les \'equations alg\'ebriques dont toutes les racines sont r\'eelles. \emph{Journal de Math\'ematiques Sp\'eciales} 4 (1895) 7.
\bibitem{Ma76} S. Mandelstam. Vortices and quark confinement in non abelian gauge theories. \emph{Physics Reports} 23 (1976) 245.
\bibitem{MO94} F. Martinelli and E. Olivieri. Approach to equilibrium of Glauber dynamics in the one phase region I and II. \emph{Communications in Mathematical Physics} 161 (1994) 447 and 487.
\bibitem{MC90} F. Masry and S. Cambanis. Trapezoidal Monte Carlo integration. \emph{SIAM Journal on Numerical Analysis} 27 (1990) 225.
\bibitem{McK81} B. D. McKay. The expected eigenvalue distribution of a large regular graph. \emph{Linear Algebra and its Applications} 40 (1981) 203.
\bibitem{MS67} H. P. McKean and I. M. Singer. Curvature and the eigenvalues of the Laplacian. \emph{Journal of Differential Geometry} 1 (1967) 43.
\bibitem{MR53} N. Metropolis, A. W. Rosenbluth, M. N. Rosenbluth, A. H. Teller and E. Teller. Equation of state calculations by fast computing machines. \emph{Journal of Chemical Physics} 21 (1953) 1087.
\bibitem{MT93} S. P. Meyn and R. L. Tweedie. \emph{Markov Chains and Stochastic Stability}. Springer, 1993.
\bibitem{Mi75} A. A. Migdal. Recursion equations in gauge field theories. \emph{Soviet Physics JETP} 42 (1975) 413.
\bibitem{DM04} P. Del Moral. \emph{Feynman Kac Formulae: Genealogical and Interacting Particle Systems with Applications}. Springer, 2004.
\bibitem{MP04} C. Morningstar and M. Peardon. Analytic smearing of SU(3) link variables in lattice QCD. \emph{Physical Review D} 69 (2004) 054501.
\bibitem{MP} C. Morningstar and M. Peardon. The glueball spectrum from an anisotropic lattice study. \emph{Physical Review D} 60 (1999) 034509.
\bibitem{MHAGH22} S. M\"uller, N. Hollmann, S. Pineda Arango, J. Grabocka and F. Hutter. Transformers can do Bayesian inference. arXiv:2112.10510, 2021 (ICLR 2022).
\bibitem{Mu81} G. M\"unster. High temperature expansions for the free energy of vortices and the string tension in lattice gauge theories. \emph{Nuclear Physics B} 180 (1981) 23.
\bibitem{Mu81b} G. M\"unster. Strong coupling expansions for the mass gap in lattice gauge theories. \emph{Nuclear Physics B} 190 (1981) 439.
\bibitem{MMGS} I. Murray, D. J. C. MacKay, Z. Ghahramani and J. Skilling. Nested sampling for Potts models. \emph{Advances in Neural Information Processing Systems} 18 (2006).
\bibitem{Ne03} R. M. Neal. Slice sampling. \emph{Annals of Statistics} 31 (2003) 705.
\bibitem{Ne73} E. Nelson. Construction of quantum fields from Markoff fields. \emph{Journal of Functional Analysis} 12 (1973) 97.
\bibitem{Ne66} E. Nelson. Derivation of the Schr\"odinger equation from Newtonian mechanics. \emph{Physical Review} 150 (1966) 1079.
\bibitem{Ne76} C. M. Newman. Fourier transforms with only real zeros. \emph{Proceedings of the American Mathematical Society} 61 (1976) 245.
\bibitem{NR94} M. A. Newton and A. E. Raftery. Approximate Bayesian inference with the weighted likelihood bootstrap. \emph{Journal of the Royal Statistical Society B} 56 (1994) 3.
\bibitem{On44} L. Onsager. Crystal statistics I. A two dimensional model with an order disorder transition. \emph{Physical Review} 65 (1944) 117.
\bibitem{OAI} OpenAI. Announcement of an existence proof for quantum Yang Mills theory, 26 October 2026; Yang Mills working group chaired by D. J. Gross, Lean code and validation assessment, 10 November 2026. \url{https://openai.com/index/building-scientific-understanding-with-academic-communities/}
\bibitem{Opp18} I. Oppenheim. Local spectral expansion approach to high dimensional expanders I: descent of spectral gaps. \emph{Discrete and Computational Geometry} 59 (2018) 293.
\bibitem{OSe} K. Osterwalder and E. Seiler. Gauge field theories on a lattice. \emph{Annals of Physics} 110 (1978) 440.
\bibitem{OS} K. Osterwalder and R. Schrader. Axioms for Euclidean Green's functions. \emph{Communications in Mathematical Physics} 31 (1973) 83 and 42 (1975) 281.
\bibitem{PPS} O. H. M. Padilla, N. G. Polson and J. G. Scott. Bayes nonparametric deconvolution with theoretical guarantees.
\bibitem{PR14} O. Papaspiliopoulos and M. Ruggiero. Optimal filtering and the dual process. \emph{Bernoulli} 20 (2014) 1999.
\bibitem{PW81} G. Parisi and Y. S. Wu. Perturbation theory without gauge fixing. \emph{Scientia Sinica} 24 (1981) 483.
\bibitem{PPR} G. Parisi, R. Petronzio and F. Rapuano. A measurement of the string tension near the continuum limit. \emph{Physics Letters B} 128 (1983) 418.
\bibitem{Pa80} G. Parisi. Recent progresses in gauge theories. In \emph{High Energy Physics 1980}, AIP Conference Proceedings 68, 1981.
\bibitem{PW27} F. Peter and H. Weyl. Die Vollst\"andigkeit der primitiven Darstellungen einer geschlossenen kontinuierlichen Gruppe. \emph{Mathematische Annalen} 97 (1927) 737.
\bibitem{Pi36} E. J. G. Pitman. Sufficient statistics and intrinsic accuracy. \emph{Mathematical Proceedings of the Cambridge Philosophical Society} 32 (1936) 567.
\bibitem{PiSh} M. K. Pitt and N. Shephard. Filtering via simulation: auxiliary particle filters. \emph{Journal of the American Statistical Association} 94 (1999) 590.
\bibitem{Po73} H. D. Politzer. Reliable perturbative results for strong interactions? \emph{Physical Review Letters} 30 (1973) 1346.
\bibitem{PR} N. G. Polson and C. P. Robert. Discussion of Vardi and Lee. \emph{Journal of the Royal Statistical Society B}, 1994.
\bibitem{PSdual} N. G. Polson and J. G. Scott. Mixtures, envelopes and hierarchical duality. \emph{Journal of the Royal Statistical Society B} 78 (2016) 701.
\bibitem{PSvl} N. G. Polson and J. G. Scott. Vertical likelihood Monte Carlo. arXiv, 2014.
\bibitem{PSX23} N. G. Polson, V. Sokolov and J. Xu. Quantum Bayesian computation. \emph{Applied Stochastic Models in Business and Industry} (2023). DOI 10.1002/asmb.2807; arXiv:2208.08068.
\bibitem{PSZ26} N. G. Polson, V. Sokolov and D. Zantedeschi. Bell's inequality, causal bounds, and quantum Bayesian computation: a unified framework. arXiv:2603.28973, 2026.
\bibitem{PS26m} N. G. Polson and V. Sokolov. The inequality (M): the core step in the Alp\"oge and Claude proof of RH. Slides, November 2026.
\bibitem{PS26t} N. G. Polson and V. Sokolov. Theta lifts and the Thorin programme. Chapter for \emph{Computing the Riemann Hypothesis}, manuscript, November 2026.
\bibitem{PS26a} N. G. Polson and V. Sokolov. \emph{Computing the Riemann Hypothesis, Volume I: Representations of Xi, Thorin Measures and the Reciprocal Xi Function}. Version 2.2.0, Zenodo, 2026. DOI 10.5281/zenodo.24086101.
\bibitem{PS26b} N. G. Polson and V. Sokolov. \emph{Computing the Riemann Hypothesis, Volume II: Cauchy Deconvolution and Certified Computation}. Version 3.0.0, Zenodo, 2026. DOI 10.5281/zenodo.24377753.
\bibitem{PS26c} N. G. Polson and V. Sokolov. \emph{Computing the Riemann Hypothesis, Volume III: Is There a Path? Routes to RH from the Thorin Programme}. Version 3.0.0, Zenodo, 2026. DOI 10.5281/zenodo.24377757.
\bibitem{PS26h} N. G. Polson and V. Sokolov. \emph{On Hilbert's 6th Problem: Probability, Mechanics and Bayesian Computation}. Version 1.1.0, independently published, 2026. DOI 10.5281/zenodo.23741447.
\bibitem{PS26l} N. G. Polson and V. Sokolov. \emph{On Langlands Conjectures: Probabilistic Methods}. Version 1.4.0, independently published, 2026. DOI 10.5281/zenodo.23904867.
\bibitem{PSgbc} N. G. Polson and V. Sokolov. Generative AI for Bayesian computation. arXiv:2305.14972, 2023.
\bibitem{PSW15} N. G. Polson, J. G. Scott and B. T. Willard. Proximal algorithms in statistics and machine learning. \emph{Statistical Science} 30 (2015) 559.
\bibitem{PSW13} N. G. Polson, J. G. Scott and J. Windle. Bayesian inference for logistic models using P\'olya Gamma latent variables. \emph{Journal of the American Statistical Association} 108 (2013) 1339.
\bibitem{Po26g} N. G. Polson. \emph{On Statistical Methods for Genetics: Fisher, Ewens and the Bayesian Programme}. Version 1.0.1, independently published, 2026. DOI 10.5281/zenodo.23509850.
\bibitem{Po96} N. G. Polson. Convergence of Markov chain Monte Carlo algorithms. In \emph{Bayesian Statistics 5}, Oxford University Press, 1996.
\bibitem{Po26b} N. G. Polson. From P\'olya's Hilfssatz II to the Thorin representation: a probabilistic synthesis. Manuscript, April 2026.
\bibitem{Po17} N. G. Polson. On Hilbert's 8th problem. arXiv:1708.02653.
\bibitem{PS14} G. P\'olya and J. Schur. \"Uber zwei Arten von Faktorenfolgen in der Theorie der algebraischen Gleichungen. \emph{Journal f\"ur die reine und angewandte Mathematik} 144 (1914) 89.
\bibitem{Po26} G. P\'olya. Bemerkung \"uber die Integraldarstellung der Riemannschen $\xi$ Funktion. \emph{Acta Mathematica} 48 (1926) 305.
\bibitem{Po77} A. M. Polyakov. Quark confinement and topology of gauge theories. \emph{Nuclear Physics B} 120 (1977) 429.
\bibitem{Pt30} E. L. Post. Generalized differentiation. \emph{Transactions of the American Mathematical Society} 32 (1930) 723.
\bibitem{Ri59} B. Riemann. Ueber die Anzahl der Primzahlen unter einer gegebenen Gr\"osse. \emph{Monatsberichte der Berliner Akademie}, 1859.
\bibitem{RR99} G. O. Roberts and J. S. Rosenthal. Convergence of slice sampler Markov chains. \emph{Journal of the Royal Statistical Society B} 61 (1999) 643.
\bibitem{RS97} G. O. Roberts and C. F. Shortland. Pricing barrier options with time dependent coefficients. \emph{Mathematical Finance} 7 (1997) 83.
\bibitem{RS95} G. O. Roberts and C. F. Shortland. The hazard rate tangent approximation for boundary hitting times. \emph{Annals of Applied Probability} 5 (1995) 446.
\bibitem{RuSe65} H. Rubin and J. Sethuraman. Probabilities of moderate deviations; Bayes risk efficiency. \emph{Sankhy\=a A} 27 (1965) 325 and 347.
\bibitem{Rub26} Claude Fable 5.2, Codex Astra 6.2 and Gemini Carbon, at the behest of M. O. Rubinstein. The Alp\"oge and Claude proof of the Riemann hypothesis explained for the Riemann zeta function. Unpublished working draft, 164 pages, 21 November 2026. Copy seen by the authors.
\bibitem{SSV11} S. Schaefer, R. Sommer and F. Virotta. Critical slowing down and error analysis in lattice QCD simulations. \emph{Nuclear Physics B} 845 (2011) 93.
\bibitem{SS05} A. D. Scott and A. D. Sokal. The repulsive lattice gas, the independent set polynomial, and the Lov\'asz local lemma. \emph{Journal of Statistical Physics} 118 (2005) 1151.
\bibitem{SW94} N. Seiberg and E. Witten. Electric magnetic duality, monopole condensation, and confinement in N=2 supersymmetric Yang Mills theory. \emph{Nuclear Physics B} 426 (1994) 19.
\bibitem{Se12} E. Seiler. Critical discussion of Tomboulis' approach to the confinement problem. \emph{Proceedings of Science} ConfinementVIII (2012) 034.
\bibitem{Se65} A. Selberg. On the estimation of Fourier coefficients of modular forms. \emph{Proceedings of Symposia in Pure Mathematics} 8 (1965) 1.
\bibitem{Se97} A. Sengupta. Gauge theory on compact surfaces. \emph{Memoirs of the American Mathematical Society} 126 (1997).
\bibitem{SZZ} H. Shen, R. Zhu and X. Zhu. A stochastic analysis approach to lattice Yang Mills at strong coupling. \emph{Communications in Mathematical Physics} 400 (2023) 805.
\bibitem{Sk06} J. Skilling. Nested sampling for general Bayesian computation. \emph{Bayesian Analysis} 1 (2006) 833.
\bibitem{SS12} A. Sly and N. Sun. The computational hardness of counting in two spin models on d regular graphs. \emph{Proceedings of the 53rd IEEE Symposium on Foundations of Computer Science} (2012) 361.
\bibitem{So01} A. D. Sokal. Bounds on the complex zeros of (di)chromatic polynomials and Potts model partition functions. \emph{Combinatorics, Probability and Computing} 10 (2001) 41.
\bibitem{So80} A. D. Sokal. An improvement of Watson's theorem on Borel summability. \emph{Journal of Mathematical Physics} 21 (1980) 261.
\bibitem{St67} F. W. Steutel. Note on the infinite divisibility of exponential mixtures. \emph{Annals of Mathematical Statistics} 38 (1967) 1303.
\bibitem{St02} G. Storvik. Particle filters for state space models with the presence of unknown static parameters. \emph{IEEE Transactions on Signal Processing} 50 (2002) 281.
\bibitem{St60} R. L. Stratonovich. Conditional Markov processes. \emph{Theory of Probability and its Applications} 5 (1960) 156.
\bibitem{SW64} R. F. Streater and A. S. Wightman. \emph{PCT, Spin and Statistics, and All That}. Benjamin, 1964 (Princeton University Press reprint, 2001).
\bibitem{SZ92} D. W. Stroock and B. Zegarlinski. The logarithmic Sobolev inequality for discrete spin systems on a lattice. \emph{Communications in Mathematical Physics} 149 (1992) 175.
\bibitem{SY82} B. Svetitsky and L. G. Yaffe. Critical behavior at finite temperature confinement transitions. \emph{Nuclear Physics B} 210 (1982) 423.
\bibitem{SW87} R. H. Swendsen and J. S. Wang. Nonuniversal critical dynamics in Monte Carlo simulations. \emph{Physical Review Letters} 58 (1987) 86.
\bibitem{Sw79} R. H. Swendsen. Monte Carlo renormalization group. \emph{Physical Review Letters} 42 (1979) 859.
\bibitem{Sz22} G. Szeg\H o. Bemerkungen zu einem Satz von J. H. Grace \"uber die Wurzeln algebraischer Gleichungen. \emph{Mathematische Zeitschrift} 13 (1922) 28.
\bibitem{Tav84} S. Tavar\'e. Line of descent and genealogical processes, and their applications in population genetics models. \emph{Theoretical Population Biology} 26 (1984) 119.
\bibitem{Te98} M. Teper. Glueball masses and other physical properties of SU(N) gauge theories in D = 3+1: a review of lattice results for theorists. arXiv:hep-th/9812187, 1998.
\bibitem{Te99} M. Teper. SU(N) gauge theories in 2+1 dimensions. \emph{Physical Review D} 59 (1999) 014512.
\bibitem{Th77} O. Thorin. On the infinite divisibility of the Pareto distribution. \emph{Scandinavian Actuarial Journal} (1977) 31.
\bibitem{Th77b} O. Thorin. On the infinite divisibility of the lognormal distribution. \emph{Scandinavian Actuarial Journal} (1977) 121.
\bibitem{To07} E. T. Tomboulis. Confinement for all values of the coupling in four dimensional SU(2) gauge theory. arXiv:0707.2179, 2007.
\bibitem{Tr91} M. D. Troutt. A theorem on the density of the density ordinate and an alternative interpretation of the Box Muller method. \emph{Statistics} 22 (1991) 463.
\bibitem{VL} Y. Vardi and D. Lee. From image deblurring to optimal investments: maximum likelihood solutions for positive linear inverse problems. \emph{Journal of the Royal Statistical Society B} 55 (1993) 569.
\bibitem{Vas17} A. Vaswani and coauthors. Attention is all you need. \emph{Advances in Neural Information Processing Systems} 30, 2017.
\bibitem{Wa22} J. L. Walsh. On the location of the roots of certain types of polynomials. \emph{Transactions of the American Mathematical Society} 24 (1922) 163.
\bibitem{WL01} F. Wang and D. P. Landau. Efficient, multiple range random walk algorithm to calculate the density of states. \emph{Physical Review Letters} 86 (2001) 2050.
\bibitem{Wat44} G. N. Watson. \emph{A Treatise on the Theory of Bessel Functions}. Second edition, Cambridge University Press, 1944.
\bibitem{We29} H. Weyl. Elektron und Gravitation I. \emph{Zeitschrift f\"ur Physik} 56 (1929) 330.
\bibitem{Wi41} D. V. Widder. \emph{The Laplace Transform}. Princeton University Press, 1941.
\bibitem{Wi74} K. G. Wilson. Confinement of quarks. \emph{Physical Review D} 10 (1974) 2445.
\bibitem{Wi98} E. Witten. Anti de Sitter space, thermal phase transition, and confinement in gauge theories. \emph{Advances in Theoretical and Mathematical Physics} 2 (1998) 505.
\bibitem{Wi91} E. Witten. On quantum gauge theories in two dimensions. \emph{Communications in Mathematical Physics} 141 (1991) 153.
\bibitem{Woit26} P. Woit. Not Even Wrong, entries of 26 October 2026 and after, with the comments of the working group quoted there. \texttt{math.columbia.edu/\~{}woit/wordpress}.
\bibitem{Ya78} S. Yakowitz, J. E. Krimmel and F. Szidarovszky. Weighted Monte Carlo integration. \emph{SIAM Journal on Numerical Analysis} 15 (1978) 1289.
\bibitem{YM54} C. N. Yang and R. L. Mills. Conservation of isotopic spin and isotopic gauge invariance. \emph{Physical Review} 96 (1954) 191.
\bibitem{Yor80} M. Yor. Loi de l'indice du lacet brownien, et distribution de Hartman Watson. \emph{Zeitschrift f\"ur Wahrscheinlichkeitstheorie} 53 (1980) 71.
\bibitem{Ki95} A. Yu. Kitaev. Quantum measurements and the abelian stabilizer problem. arXiv:quant-ph/9511026, 1995.
\bibitem{MM79} Yu. M. Makeenko and A. A. Migdal. Exact equation for the loop average in multicolor QCD. \emph{Physics Letters B} 88 (1979) 135.
\bibitem{Za69} M. Zakai. On the optimal filtering of diffusion processes. \emph{Zeitschrift f\"ur Wahrscheinlichkeitstheorie} 11 (1969) 230.
\bibitem{Zi96} M. R. Zirnbauer. Supersymmetry for systems with unitary disorder: circular ensembles. \emph{Journal of Physics A} 29 (1996) 7113.
\end{thebibliography}
